\section{The base theta map and its finite stabilizer}
\label{sec:base-theta-v032}

Fix a layer $i$ and the base tuple $t_0$. The theta integral is first
defined on vectors in the lowest weight space. We need its values in
two forms. As automorphic functions they give complex periods. As
coordinates on the rational model they enter the integral matrix whose
minors we estimate. To compare the two, we keep the chosen Weil action,
the embeddings of the coefficient field, the measures and the basis of
the target fixed throughout.

% excerpt: Here is the notation used throughout.  Put || % END BTH15

\subsection*{From the integral to the base rows}

Convergence gives a value of the theta integral on each vector of the
chosen basis of $H_i$. To obtain a vector of $W_i$, those values must
extend to an element of $\mathrm{Wall}_i$, the multiplicity space in
the specified realization of $W_i$. The extension uses smoothness,
moderate growth, the vanishing of transvection constant terms, the
compact-type equivariance and annihilation by the lowering operators.
These conditions identify the cuspidal automorphic realization in
which the coordinates of the period will be taken. The two steps are
the \href{https://github.com/CocoSgt/riemann-hypothesis-lean/blob/7cbf7a0c32fca1badf05931879a37c3f8a2b03b4/src/Development/AbelianVarietyTheory/Proofs/IwasawaDecomposition/ExistsThetaLiftLinearMapCuspFormMultiplicity.lean#L180-L224}{pointwise extension}
and its \href{https://github.com/CocoSgt/riemann-hypothesis-lean/blob/7cbf7a0c32fca1badf05931879a37c3f8a2b03b4/src/Development/AbelianVarietyTheory/Proofs/IwasawaDecomposition/ExistsThetaLiftLinearMapCuspFormMultiplicity.lean#L77-L158}{realization in the multiplicity space}.

There is also a question about the field of definition. Suppose a
finite set of linear functionals separates the vectors in the complex
target. If their values on the chosen target basis and on the theta
images of the rational source vectors all belong to the same embedded
coefficient field, solving the resulting linear equations recovers the
theta coordinates in that field. The functionals must test the two
sets of vectors in the same realization. The
\hyperref[ms-manuscript-notes-theta-selected-theta-rational-tests-md]{rationality argument}
identifies such a set on a selected bundle.

Finally the rows used in the period estimates must be the rows of this
matrix, with its prescribed adjoint. The
\hyperref[ms-manuscript-notes-construction-selected-base-rows-md]{base-row comparison}
uses the identification of the geometric and analytic readings at the
selected tuple, including its phase and basis, to make that last
identification. These conclusions apply once the bundle and its two
realizations have been supplied. Choosing such bundles on a whole
arithmetic window is the separate construction of
\hyperref[ms-manuscript-notes-arithmetic-common-design-window-transfer-md]{the common family}.
